\section{Approach}\label{sec:approach}

\paragraph*{Vocabulary}
The domain declares the policy-visible variables, the instance constants and the features in a small specification
next to its PRISM model. For gridworld (Section~\ref{sec:evaluation}) the constants are the grid size and the goals'
coordinates; the features say whether a move hits a wall, enters the moving obstacle's patrol path, enters a goal
that is not due yet, could slip onto one of those, goes towards the next goal, ends where the obstacle will be next,
or (optionally) shortens the path to the next goal. Each feature is an expression in the rule language, so the rules
compile to a PRISM module that synchronises on the policy actions, and checking a rule list is ordinary model
checking of the composed model.

\paragraph*{Training on several instances}
A training set is a group of small instances. Each round, the LLM sees every member (description, map, constant
values), the vocabulary and the requirements, and returns one rule list as JSON (constrained decoding bounds its
size). The list is parsed for each member, with that member's constants, and verified on each. The round's
candidate is scored by the members' worst-case failures and shortfalls, summed, and the best candidate so far is
kept. The set is solved when every member passes in the worst case.

\paragraph*{Feedback}
If the kept rules fail on some member, the next prompt reports, for the member that fails worst, the best and worst
case of every requirement and the rules' coverage. If no completion of the rules can pass (some best case fails, or a
multi-objective query shows that no single scheduler meets all thresholds), the LLM is asked to rewrite the list and
is shown how much of the lost probability each rule's states carry: expected visits within a horizon times the
probability still at stake there. Otherwise only open choices are to blame, and it is asked for additional rules,
shown the uncovered situations where the worst case loses the most. Feedback locates where probability is lost; it
never reveals the optimal action.

\subsection{Refinement by the verifier}\label{sec:approach-refine}
LLMs write \any{} rules readily: they state which moves are acceptable without committing to one. In the worst case,
an adversary chooses among those moves, and the loop optimises the rules against that adversary. Such rules can be
safe on the training instances only because they forbid too much, and they rarely carry over.

We therefore let the verifier resolve the choices. For an order $\pi$ of the actions, $\mathit{fix}_\pi(R)$ replaces
each rule $(c, \any)$ by the rules $(c, \pi_1), \dots, (c, \pi_{|A|})$: a state the rule decides takes the first
action in $\pi$ that the rule allows. Each round, after the LLM answers, the loop verifies $\mathit{fix}_\pi(R)$ on
the training members for every order $\pi$ ($|A|! = 24$ in our domains) and keeps the one that scores best. From then
on the round works with $\mathit{fix}_\pi(R)$: it is what is verified, kept, shown to the LLM (with one rule per
action) and, at the end, frozen. The LLM is told that this happens. The order is chosen on the training instances
only, so the frozen rules are certified on held-out instances exactly as before.

\paragraph*{Certification on held-out instances}
After the last round, the kept rule list is frozen. For every held-out instance, it is parsed with that instance's
constants and certified with PRISM; worst cases close to a threshold are re-checked with interval iteration so that
verdicts are exact. A rule list that does not parse on an instance (it cannot, in general mode, unless it uses a
constant the instance lacks) would count as a failure.
